% p.7 — continuation of الباب الثاني (columns 11–15 and خلاصة)
% Layout note: the lithograph has three column-groups of unequal structure:
%  (a) right: اسم ظرف زمان ومكان (11) / مثال / تمرين, then اسم آلة (12) in three variants, then اسم مرة (14);
%  (b) middle: اسم تفضيل (13) / مثال / تمرين, and اسم مبالغة (15) / مثال / تمرين;
%  (c) left: خلاصة — the 15 patterns with الخ / ditto marks, then تمرين.
% They are set side by side (right to left) with the tables environment, as in the lithograph.
\mspage{7}
\begin{tables}
\begin{sarf}{|c|c|c|}
\thead{(11) اسم ظرف زمان ومكان} & \thead{مثال} & \thead{تمرين}\\\hline
مَفْعَلٌ & مَفْتَحٌ & مَطْبَخٌ\\
مَفْعَلَانِ & مَفْتَحَانِ & مَطْحَنٌ\\
مَفَاعِلُ & مَفَاتِحُ & مَقْتَلٌ\\\hline
\thead{(12) اسم آلة} & \thead{مثال} & \thead{تمرين}\\\hline
مِفْعَلٌ & مِفْتَحٌ & مِطْحَنٌ\\
مِفْعَلَانِ & مِفْتَحَانِ & مِسْمَعٌ\\
مَفَاعِلُ & مَفَاتِحُ & مِنْظَرٌ\\\hline
\thead{(12) اسم آلة} & \thead{مثال} & \thead{تمرين}\\\hline
مِفْعَلَةٌ & مِفْتَحَةٌ & مِدْفَعَةٌ\\
مِفْعَلَتَانِ & مِفْتَحَتَانِ & مِضْرَبَةٌ\\
مَفَاعِلُ & مَفَاتِحُ & مِنْصَرَةٌ\\\hline
\thead{(12) اسم آلة} & \thead{مثال} & \thead{تمرين}\\\hline
مِفْعَالٌ & مِفْتَاحٌ & مِرْفَاعٌ\\
مِفْعَالَانِ & مِفْتَاحَانِ & مِقْطَاعٌ\\
مَفَاعِيْلُ & مَفَاتِيْحُ & مِرْوَاحٌ\\\hline
\thead{(14) اسم مرة} & \thead{مثال} & \thead{تمرين}\\\hline
فَعْلَةٌ & فَتْحَةٌ & نَصْرَةٌ\\
فَعْلَتَانِ & فَتْحَتَانِ & ضَرْبَةٌ\\
فَعَلَاتٌ & فَتَحَاتٌ & مَرَّةٌ\\
\end{sarf}
\begin{sarf}{|c|c|}
\thead{(13) اسم تفضيل} & \thead{(15) اسم مبالغة}\\\hline
أَفْعَلُ & فَعَّالٌ\\
أَفْعَلَانِ & فَعَّالَانِ\\
أَفْعَلُوْنَ & فَعَّالُوْنَ\\
أَفَاعِلُ & فَعَّالَةٌ\\
فُعْلٰى & فَعَّالَتَانِ\\
فُعْلَيَانِ & فَعَّالَاتٌ\\\cline{2-2}
فُعْلَيَاتٌ & \thead{مثال}\\\cline{2-2}
فُعَلٌ & فَتَّاحٌ\\\cline{1-1}
\thead{مثال} & فَتَّاحَانِ\\\cline{1-1}
أَفْتَحُ & فَتَّاحُوْنَ\\
أَفْتَحَانِ & فَتَّاحَةٌ\\
أَفْتَحُوْنَ & فَتَّاحَتَانِ\\
أَفَاتِحُ & فَتَّاحَاتٌ\\\cline{2-2}
فُتْحٰى & \thead{تمرين}\\\cline{2-2}
فُتْحَيَانِ & قَطَّاعٌ\\
فُتْحَيَاتٌ & ضَرَّابٌ\\
فُتَحُ & قَتَّالٌ\\\cline{1-1}
\thead{تمرين} & طَحَّانٌ\\\cline{1-1}
أَنْصَرُ & طَبَّاخٌ\\
\end{sarf}
\begin{sarf}{|c|c|}
\multicolumn{2}{|c|}{\thead{خلاصة}}\\\hline
فَعَلَ & \ilkh\\
يَفْعَلُ & \ilkh\\
فَاعِلٌ & \ditto\\
فُعِلَ & \ditto\\
يُفْعَلُ & \ditto\\
مَفْعُوْلٌ & \ditto\\
اِفْعَلْ & \ditto\\
لَا تَفْعَلْ & \ditto\\
مَفْعَلٌ & \ditto\\
مِفْعَلٌ & \ditto\\
مِفْعَلَةٌ & \ditto\\
مِفْعَالٌ & \ditto\\
أَفْعَلُ & \ditto\\
فَعْلَةٌ & \ditto\\
فَعَّالٌ & \ditto\\\hline
\multicolumn{2}{|c|}{\thead{تمرين}}\\\hline
\multicolumn{2}{|c|}{طَحَنَ}\\
\multicolumn{2}{|c|}{دَفَعَ}\\
\multicolumn{2}{|c|}{قَطَعَ}\\
\end{sarf}
\end{tables}
